% Options for packages loaded elsewhere
\PassOptionsToPackage{unicode}{hyperref}
\PassOptionsToPackage{hyphens}{url}
\PassOptionsToPackage{dvipsnames,svgnames,x11names}{xcolor}
\documentclass[
  11pt,
  a4paper]{article}
\usepackage{xcolor}
\usepackage[margin=1.6cm]{geometry}
\usepackage{amsmath,amssymb}
\setcounter{secnumdepth}{-\maxdimen} % remove section numbering
\usepackage{iftex}
\ifPDFTeX
  \usepackage[T1]{fontenc}
  \usepackage[utf8]{inputenc}
  \usepackage{textcomp} % provide euro and other symbols
\else % if luatex or xetex
  \usepackage{unicode-math} % this also loads fontspec
  \defaultfontfeatures{Scale=MatchLowercase}
  \defaultfontfeatures[\rmfamily]{Ligatures=TeX,Scale=1}
\fi
\usepackage{lmodern}
\ifPDFTeX\else
  % xetex/luatex font selection
\fi
% Use upquote if available, for straight quotes in verbatim environments
\IfFileExists{upquote.sty}{\usepackage{upquote}}{}
\IfFileExists{microtype.sty}{% use microtype if available
  \usepackage[]{microtype}
  \UseMicrotypeSet[protrusion]{basicmath} % disable protrusion for tt fonts
}{}
\makeatletter
\@ifundefined{KOMAClassName}{% if non-KOMA class
  \IfFileExists{parskip.sty}{%
    \usepackage{parskip}
  }{% else
    \setlength{\parindent}{0pt}
    \setlength{\parskip}{6pt plus 2pt minus 1pt}}
}{% if KOMA class
  \KOMAoptions{parskip=half}}
\makeatother
\setlength{\emergencystretch}{3em} % prevent overfull lines
\providecommand{\tightlist}{%
  \setlength{\itemsep}{0pt}\setlength{\parskip}{0pt}}
\definecolor{accent}{HTML}{2E6E8E}
\usepackage{fontawesome5}
\usepackage{titlesec}
\usepackage{enumitem}
\setlist[itemize]{leftmargin=1.2em,topsep=4pt,itemsep=2pt,parsep=0pt}
\titleformat{\section}{\large\bfseries\color{accent}}{}{0em}{}[\vspace{2pt}{\titlerule[0.4pt]}]
\titlespacing{\section}{0pt}{12pt}{5pt}
\titleformat{\subsection}{\large\bfseries\color{accent}}{}{0em}{}[\vspace{2pt}{\titlerule[0.4pt]}]
\titlespacing{\subsection}{0pt}{12pt}{5pt}
\pagestyle{empty}
\usepackage{bookmark}
\IfFileExists{xurl.sty}{\usepackage{xurl}}{} % add URL line breaks if available
\urlstyle{same}
% fallback for those not using the hyperref driver hyperxmp:
\makeatletter
\@ifundefined{xmpquote}{\newcommand{\xmpquote}[1]{#1}}{}
\makeatother
\hypersetup{
  colorlinks=true,
  linkcolor={accent},
  filecolor={Maroon},
  citecolor={Blue},
  urlcolor={accent},
  pdfcreator={LaTeX via pandoc}}

\author{}
\date{}

\begin{document}

\begin{center}
{\Huge\bfseries\color{accent} Oriol Vilaró Roca\par}
\vspace{2pt}
{\large Desenvolupador de videojocs i d'aplicacions mòbils\par}
\vspace{6pt}
{\small
\faIcon{phone}\ +34 625 91 82 79 \quad
\faIcon{envelope}\ \href{mailto:ovilaro@fastmail.com}{ovilaro@fastmail.com}\par
\faIcon{map-marker}\ Raval Josep Pané, 16 · 25211 Tarroja de Segarra, Lleida, Espanya\par}
\end{center}
\vspace{8pt}

\subsection{Perfil}\label{perfil}

Desenvolupador de videojocs i d'aplicacions mòbils; autodidacte,
m'agrada aprendre nous llenguatges i tecnologies. Tinc experiència
desenvolupant en solitari i en equip, emprant eines per a gestionar les
tasques del projecte, la documentació i el codi font.

\subsection{Experiència}\label{experiuxe8ncia}

\textbf{Desenvolupador de videojocs · DABADU GAMES} --- 2017--2024

\begin{itemize}
\tightlist
\item
  Desenvolupament en Unity3D de tres videojocs per a iOS i Android,
  destacats en App Store i Google Play en múltiples ocasions.
\item
  Desenvolupador líder d'un videojoc amb la llicència d'una popular
  franquícia japonesa, publicat per Netflix en iOS i Android i portat a
  Nintendo Switch.
\item
  Desenvolupament en Godot del prototip d'un projecte més ambiciós, amb
  màquines d'estat i generació de nivells procedurals.
\end{itemize}

\textbf{Cap d'equip i desenvolupador d'aplicacions mòbils · MODDITY} ---
2014--2017

\begin{itemize}
\tightlist
\item
  Gestió del desenvolupament d'aplicacions per a dispositius iOS i
  Android, dirigint equips de programadors.
\item
  Especialització en desenvolupament per a iOS.
\end{itemize}

\textbf{Programador d'aplicacions mòbils · BAZINGA SYSTEMS} ---
2011--2014

\begin{itemize}
\tightlist
\item
  Programació d'aplicacions iOS per a iPhone i iPad.
\end{itemize}

\textbf{Programador d'aplicacions de BBDD · BONÀREA} --- 2005--2010

\begin{itemize}
\tightlist
\item
  Programació i gestió del control d'un magatzem de productes
  alimentaris: gestió d'estoc, preparació de comandes, entrades i
  sortides, i gestió horària del personal.
\end{itemize}

\subsection{Educació}\label{educaciuxf3}

\textbf{CFGS Desenvolupament d'Aplicacions Informàtiques · La Salle
Gràcia, Barcelona} --- 2002--2004

\textbf{Programació en C++ · CIPSA} --- 2001--2002

\subsection{Certificacions}\label{certificacions}

\textbf{Desenvolupament d'aplicacions mòbils per a iOS i Android amb
Flutter} --- CIFO L'Hospitalet (SOC) · 2025 240 h · Apte (Notable) ·
Codi: IFCD0002-1

\subsection{Aptituds}\label{aptituds}

\begin{itemize}
\tightlist
\item
  Autodidacte
\item
  Treball en equip
\item
  Agile development
\end{itemize}

\subsection{Idiomes}\label{idiomes}

\begin{itemize}
\tightlist
\item
  Català · nivell alt
\item
  Espanyol · nivell alt
\item
  Anglès · nivell professional bàsic
\end{itemize}

\end{document}
